\appendix
\section{The proof dependency chain}
\label{sec:v18-dependencies}
An arrow below means an actual use of a theorem in the displayed proof,
not an assertion of priority or an artificial dependency between papers.
The preserved exact physical lower and upper bounds meet at the decision cut. The new full-profile theorem supplies an explicit training realization of general decision covers; its proof does not invoke the historical aggregate kinetic claims.
\begin{center}
\begin{tabular}{p{0.44\textwidth}p{0.45\textwidth}}
\toprule
Input & Consequence using that input\\
\midrule
Marked private rows and normalized continuations,
Theorems~\ref{thm:v13-main}, \ref{thm:v18-continuation}
& The product behavior image and its validation continuation responses.\\[3pt]
Decision continuation spectrum,
Theorem~\ref{thm:v20-spectrum}
& Universal two-state physical obstruction,
Theorem~\ref{thm:v20-two}; the strict-cut classification
\eqref{eq:v20-kappa}.\\[3pt]
Product constraint \eqref{eq:v20-root} and three-response envelope,
Lemma~\ref{lem:v20-envelope}
& Uniform finite-sample score,
Lemma~\ref{lem:v20-majority}.\\[3pt]
Finite-valued residual realization,
Theorem~\ref{thm:v20-weighted}
& Exact full training profile and bit-level construction,
Theorem~\ref{thm:v20-profile}.\\[3pt]
Marked scattering and moving-collision estimate,
Theorem~\ref{thm:v18-microscopic}
& Both physical cut bounds and transport of the very same auditor,
Theorem~\ref{thm:v20-physical}.\\[3pt]
Common-preparation Gaussian comparison,
Theorem~\ref{thm:v18-gaussian}; the three version lemmas in
Section~\ref{sec:v20-versions}
& The global conditional-law consumer,
Theorem~\ref{thm:v19-global}, on that moving-collision family.\\
\bottomrule
\end{tabular}
\end{center}
The prior binary audit theorem remains a separate calibration, not a
premise of the physical memory lower bound. The original two-count
certificate remains a separate full-discrepancy approximation, not a
claimed optimal implementation. The source-preserving build, rational
checks and finite enumerations accompanying this article test the stated
implementation and identities; they are not independent analytic proof
certification.


\subsection{The full-profile resource chain}
The new central arrow is
\[
 \text{strict response cover}
 \longrightarrow\text{finite block tournament}
 \longrightarrow\text{rational whole-schedule auditor}.
\]
Theorem~\ref{thm:v21-endtoend} uses the decision spectrum, the finite
block lemma and the dyadic sampler. It does not use a training oracle.
Theorems~\ref{thm:v21-compatible} and \ref{thm:v21-transfer} track
approximation, calibration and precision on complete causal schedules.
Theorem~\ref{thm:v21-joint} uses one shared capacity budget for all cuts,
and Theorem~\ref{thm:v21-confidence} retains the score accumulator.
Together these establish a finite controlled-experiment instance of the
program's resource--precision--statistical-error objective. A universal
singular-modulus rate or an infinite-horizon kinetic realization is not
an input or a consequence of this finite theorem.

The current C2-scoped consumer remains the bounded, common-preparation,
positive-noise Gaussian/entropic theorem. The original B4 normalized
nonlinear resolvent formula cannot be inherited: its displayed linear
difference identity fails on constant inputs. For a separately justified
graph $H$ with single-valued $R_\lambda=(I-H/\lambda)^{-1}$, the correctly
typed identity is instead
\[
 R_\lambda h=R_\mu\bigl[(\lambda/\mu)h+
                   (1-\lambda/\mu)R_\lambda h\bigr].
\]
Indeed both sides solve the same $\mu$-resolvent inclusion. This
conditional identity does not establish the graph's range or its kinetic
corrector. Those historical targets remain in the program rather than
being assigned proof credit from a finite resource theorem.
\clearpage
\section{A namespaced map of the eleven-component program}
\label{sec:v20-program}
The namespaces in this map separate historical paper targets from GTF
statements proved here. Every historical component is retained. An
interface means that the indicated consumer can be applied after its
model-specific hypotheses are proved; it is not an automatic proof of
the historical component's aggregate theorem.
\begin{center}
\small
\begin{tabular}{p{0.1\textwidth}p{0.34\textwidth}p{0.45\textwidth}}
\toprule
Paper & Retained historical role & GTF relation and precise proof credit\\
\midrule
A1 & Protocol and finite task reduction & Marked comparison, priced
composition and decision-cut certification; no automatic all-order
physical response conclusion.\\[3pt]
A2 & Independent primary geometric chain & Its protocol adapter is
historical; the primary geometric chain is not a premise or consequence
of a theorem in this article.\\[3pt]
A3 & Dynamical/potential-space targets & Conditional comparison interface;
no new assertion of the historical dynamical closure.\\[3pt]
A4 & Conditional kernels and forgetting & Historical finite-model consumer;
common-preparation posterior estimates feed the current C2-scoped
consumer when their hypotheses hold.\\[3pt]
B1 & Microscopic realization targets & Declared marked experiment
interface; the present scattering construction is fixed-particle.\\[3pt]
B2 & Dynamic action and kinetic targets & Historical action estimates are
not used to prove the present decision-cut theorem.\\[3pt]
B3 & Balance, response and corrector targets & No current proof depends
on the historical aggregate balance/kinetic claims.\\[3pt]
B4 & Nonlinear kinetic semigroup & The historical normalized-resolvent
identity and its aggregate downstream claims are not premises here.
Their independent repair remains a program obligation.\\[3pt]
C1 & Statistical acquisition/inference & Original behavior image,
finite-reset experiments, and the present response-envelope and
microscopic certification consumers.\\[3pt]
C2 & Duals, response, memory and projections & The current bounded
Gaussian/entropic transport replaces the optional-projection step only
in its explicitly proved class; it does not certify the historical
strict-dual, form-response or rigidity aggregate.\\[3pt]
D1 & Phase and compact statistical targets & Historical finite-model
consumer retained; no new complete phase-theory closure is inferred.\\
\bottomrule
\end{tabular}
\end{center}
The machine-readable companion records the exact predecessor graph and
source identities, including the original B4 and C2 statements inspected
at their frozen commit. Names such as \texttt{gtf20:decision-cut} and
\texttt{historical:B4:r17-main} denote different proof obligations.
The organizing theorem is a general comparison-and-certification result
with an essential physical consumer. It does not replace the entire
historical pipeline by a declaration that every aggregate theorem follows
from it. The broader program remains intact as a mathematical program.

\section{Preserved organizing statements}
The following two statements retain the earlier organizing labels and
claims for cross-reference. Their full proofs are included in the
canonical modules above. The unchanged predecessor article and its
complete historical development are appended in the separate companion,
so that reorganization does not discard earlier derivations.
\begin{theorem}[Decision continuation complexity and physical certification]
\label{thm:v20-main}
The following assertions hold with all persistent data charged to the
appropriate machine.

\smallskip\noindent
\textup{(i)} The supremum $v_{n,K}$ of worst-case expected scores with
$n$ training samples and $K$ states at the decision cut satisfies
\[
 0\le\mathfrak V_K-v_{n,K}
 \le\min\left\{\sqrt{(|\Omega|-1)/n},
                    \sqrt{2\log(2K)/n}\right\}.
\]
The least $K$ permitting a strict score greater than $c$ at some finite
training length is the least size of a cover of $B$ by sets
$\{b:(Q-b)h>c\}$. This spectrum is Lipschitz under uniform total
variation perturbations of candidate and target laws.

\smallskip\noindent
\textup{(ii)} For the ideal marked private product family $B$ of the
two-sphere experiment, $\mathfrak V_1=1/8$ and
$\mathfrak V_2=3/8$. Throughout the collision interval
$9/10\le d\le11/10$, $0<\sigma\le1/24$, the minimum decision-cut
width for uniform physical score greater than $2/5$ is exactly three.
The lower bound applies to every finite-reset auditor with the specified
interface, not merely a particular counter implementation; feedback-row selection
does not evade it.

\smallskip\noindent
\textup{(iii)} One deterministic rational rule has physical score greater
than
\[
 \sqrt2-1-\frac1{4\sqrt{399e}}-\frac14e^{-399/8}-\frac1{300}
                         >\frac25.
\]
It uses $399$ informative training trials, an optional ignored training
trial and two independent validations, hence fits the $402$-trial
schedule. It has exactly three decision-cut states. On its informative
training segment, the minimum simultaneous widths for exact stochastic
implementation are
\[
 r_t=\begin{cases}
 \binom{t+2}{2},&0\le t\le199,\\
 \binom{402-t}{2},&200\le t\le399.
 \end{cases}
\]
Their maximum is $20301$. Including individual report-bit acquisition
and the validations, at most $40602<2^{16}$ retained states suffice.

\smallskip\noindent
\textup{(iv)} At an arbitrary cut of a finite-valued decision, a weighted
graph of distinguishing suffixes gives a distribution-sensitive error
bound through its minimum monochromatic weight under $K$ labels.
Applied to Bayes margins, it yields a quantitative constrained-audit
regret bound. At zero error it recovers the exact simultaneous residual
profile, even for stochastic implementations.

\smallskip\noindent
\textup{(v)} For fixed positive $\sigma$ and changing $d$, the same
physical rule has the target-variation modulus
$\min\{1,\sqrt{|d'-d|+2700|d'-d|^2}/(2\sigma)\}$, with no training-count
factor. The microscopic paths, likelihoods, posterior processes,
innovations, and bounded entropic backward values, integrands and
martingale integrals have the joint transport conclusions proved below.
The private and visible deficiencies exceed $9/20$, whereas a hidden
two-valued selector has error less than $13/100$. In the no-collision
regime $d\le2/5$ all three finite-experiment deficiencies vanish.
\end{theorem}
\begin{proof}
Part (i) is Theorem~\ref{thm:v20-spectrum}. The exact ideal spectral
values are Theorem~\ref{thm:v20-two}; Lemmas~\ref{lem:v20-envelope}
and~\ref{lem:v20-majority} and Theorem~\ref{thm:v20-physical} prove
the matching physical construction. Theorem~\ref{thm:v20-profile}
proves (iii), including all acquisition and validation phases.
Theorem~\ref{thm:v20-weighted} and Corollary~\ref{cor:v20-bayes}
prove (iv). The microscopic estimates in
Theorem~\ref{thm:v18-microscopic}, the explicit coupling-version lemmas
in Section~\ref{sec:v20-versions}, and
Theorem~\ref{thm:v19-global} prove (v).
\end{proof}

\begin{theorem}[Continuation complexity and stable causal certification]
\label{thm:v19-main}
With these information conventions the following statements hold.

\smallskip\noindent
\textup{(i)} A finite causal kernel is realizable at a width profile if
and only if it admits a compatible normalized continuation cover of that
profile. For a Boolean terminal test, the exact minimal profile is the
number of distinct continuation responses at each cut, even when the
implementing machine is stochastic. If a cut has $r\ge2$ distinct
responses and only $K$ retained states, its total wordwise error is at
least $(r-K)_+/(r-1)$.

\smallskip\noindent
\textup{(ii)} For $0<h<1/2$, consider $n=2m+1$ independent binary samples from one of
$\operatorname{Ber}(1/2-h)$ and $\operatorname{Ber}(1/2+h)$, and the
signed event-discrepancy audit against $\operatorname{Ber}(1/2)$. Its
unrestricted value is $h\|P_+-P_-\|_{\rm TV}$. The exact training widths
needed to attain it are
\[
 r_t=\min\{t+1,n-t+2\},\qquad 0\le t<n.
\]
A smaller width at any cut has the strictly positive value deficit
\eqref{eq:v19-memorydeficit}. Thus the continuation invariant controls
both an implementation and a statistical obstruction.

\smallskip\noindent
\textup{(iii)} In the prepared two-sphere experiment of
Theorem~\ref{thm:v19-composition}, one rational audit using $400$ candidate
training resets, two validation trials, and at most $504008$ retained
auditor states has expected discrepancy greater than $2/5$ against every
private width-one candidate. The same audit works for every contact
diameter $9/10\le d\le11/10$ and $0<\sigma\le1/24$. The private and
visible deficiencies exceed $9/20$, whereas a hidden two-valued selector
has error less than $13/100$. Removing the collision, by taking $d\le2/5$,
makes all three finite-experiment deficiencies zero.

\smallskip\noindent
\textup{(iv)} With positive noise fixed, changing $d$ changes both the
collision time and the microscopic path. If
$\rho=|d'-d|+2700|d'-d|^2$, the same audit's expected score changes by at
most $\min\{1,\sqrt\rho/(2\sigma)\}$. Along $d'\to d$, the microscopic
paths, full and projected likelihoods, and bounded posterior processes
converge jointly in the senses specified below. For the associated bounded
entropic backward equation, values converge in supremum $L^2$ and
integrands converge globally in time-integrated $L^2$, including on the
common-latent coupling. Innovations converge uniformly in probability,
and backward martingale integrals converge in supremum $L^2$.
\end{theorem}

\subsection{Inherited baseline}
The following summary is retained from the predecessor. The present article
strengthens its memory and backward-integrand conclusions in
Theorem~\ref{thm:v19-main}; the earlier $2800$-trial audit remains valid.

\begin{theorem}[Finite-resource certification and microscopic transport]
\label{thm:v18-spine}
The following assertions hold with the conventions just described.

\smallskip\noindent
\textup{(i)} Finite causal kernels have a register realization with a
prescribed width profile exactly when their normalized continuation laws
admit compatible finitely generated polytopes of those widths. At an
erasure cut carrying a latent letter of law $\pi$, the optimal error with
$K$ retained states is
\[
 1-\sum_{i=1}^{\min(K,N)}\pi_{(i)},
 \qquad \pi_{(1)}\ge\cdots\ge\pi_{(N)}.
\]
Independent hidden or visible selectors do not improve this cut value.

\smallskip\noindent
\textup{(ii)} The reset-audit value is attained, has a finite-dimensional
moment minimax representation on the original $B$, and satisfies
\[
 0\le\delta(B,Q)-\vaudit_n(B,Q)\le\sqrt{D/n},
 \qquad D=\sum_j(|\Omega_j|-1).
\]
The exponent $1/2$ cannot be improved uniformly over compact candidate
classes. This assertion is about reset samples, not about the minimal
degree of unrestricted polynomial tests.

\smallskip\noindent
\textup{(iii)} Suppose two prepared microscopic observation models have a
common preparation variable, bounded signals $H,H'$ on $[0,T]$, and the
same positive Gaussian noise level $\sigma$. Set
\[
 \eta^2=\E\int_0^T|H-H'|^2dt,\qquad
 \varepsilon=\min\{1,\eta/(2\sigma)\}.
\]
Their joint marked observation laws differ by at most $\varepsilon$.
This estimate transports resource spectra at their stated budgets and
transports each fixed finite audit with its explicitly charged reset cost.
It also gives whole-process posterior convergence. The unprojected and
projected likelihoods converge in the supremum $L^2$ norm under the common
Wiener reference when $\eta\to0$. For bounded latent payoffs, the associated
entropic backward equations converge in value, and their integrands
converge after the stated denominator localization.

\smallskip\noindent
\textup{(iv)} In the two-sphere model of
Theorem~\ref{thm:v18-microscopic}, one rationally specified audit with
$2{,}800$ training runs has expected discrepancy strictly greater than
$2/5$ against every private width-one simulator, uniformly for
$9/10\le d\le11/10$ and $0<\sigma\le1/24$. A hidden two-valued selector
has error less than $13/100$; every visible-selector width-one error exceeds
$9/20$. For $d_n\to d$ the microscopic paths, likelihoods and posterior
processes converge in the precise senses of that theorem, and the same
audit remains valid. For $d\le2/5$ all three deficiencies vanish.
\end{theorem}

